\subsection{Первый полет: Взлет и посадка}

Управление дроном состоит из нескольких простых этапов. Рассмотрим их по порядку.

\subsubsection*{1. Запуск моторов (Arming)}
Первым делом нужно <<завести>> дрон. Команда \texttt{arm()} запускает моторы на холостых оборотах.

\begin{codefile}[python]{python/examples/02_pioneer_sdk/09_01_arm_disarm.py}
\end{codefile}

\begin{warnbox}[Безопасность]
Всегда используйте блок \texttt{try...finally}, чтобы гарантировать остановку моторов (\texttt{disarm}) даже в случае ошибки в программе или нажатия \texttt{Ctrl+C}.
\end{warnbox}

\subsubsection*{2. Взлет и посадка}
Теперь, когда моторы работают, можно взлететь. Команда \texttt{takeoff()} поднимает дрон на высоту около 1 метра.

\begin{codefile}[python]{python/examples/02_pioneer_sdk/09_02_takeoff_land.py}
\end{codefile}

Обратите внимание:
\begin{enumerate}
    \item \texttt{takeoff()} — автоматический взлет.
    \item \texttt{time.sleep(5)} — ожидание (висение) в течение 5 секунд.
    \item \texttt{land()} — автоматическая посадка.
\end{enumerate}
